\expandafter\def\csname HistoricalNoteText3\endcsname{Der analytischen Vollständigkeit halber merke ich an, dass sich der Himmelskörper in diesem Fall am
$\left\{\begin{array}{c}\AscendingNode\\\DescendingNode\end{array}\right\}$ befindet, je nachdem, ob $\sin(\odot-\alpha)$ und $\sin(\alpha-\AscendingNode)$
$\left\{\begin{array}{l}\text{gleiche}\\\text{entgegengesetzte}\end{array}\right\}$ Vorzeichen haben.}
\clearpage\SourceLeaf{92}%
\pdfbookmark[1]{Vorschriften zur Bestimmung heliozentrischer Positionen}{vorschriften1802}
\begin{center}
{\small VORSCHRIFTEN, WIE MAN AUS DER GEOZENTRISCHEN LÄNGE UND BREITE EINES HIMMELSKÖRPERS,
DER POSITION SEINES KNOTENS, DER NEIGUNG SEINER BAHN, DER LÄNGE DER SONNE UND IHREM
ABSTAND VON DER ERDE DIE HELIOZENTRISCHE LÄNGE DES HIMMELSKÖRPERS IN SEINER BAHN,
SEINEN WAHREN ABSTAND VON DER SONNE UND SEINEN WAHREN ABSTAND VON DER ERDE
ABLEITET.\par}
\medskip\hrule\medskip
{\small Monatliche Correspondenz zur Beförderung der Erd- und Himmels-Kunde,\par
herausgegeben vom Freiherrn \textsc{von Zach}. Juni 1802.\par}
\medskip\hrule
\end{center}\medskip
\begin{center}\textit{Bedeutung der Zeichen:}\end{center}
\noindent\begin{minipage}[t]{.48\textwidth}\small
\begin{center}\textit{Gegeben;}\end{center}
\begin{tabular}{@{}lp{.78\linewidth}@{}}
$\AscendingNode$ & \raggedright Länge des aufsteigenden Knotens.\tabularnewline
$\odot$ & \raggedright Länge der Sonne.\tabularnewline
$\alpha$ & \raggedright Geozentrische Länge des Himmelskörpers.\tabularnewline
$\ObservedBeta$ & \raggedright Geozentrische Breite.\tabularnewline
$i$ & \raggedright Neigung der Bahn.\tabularnewline
$R$ & \raggedright Abstand der Sonne von der Erde.
\end{tabular}
\end{minipage}\hfill
\begin{minipage}[t]{.48\textwidth}\small
\begin{center}\textit{Gesucht:}\end{center}
\begin{tabular}{@{}lp{.78\linewidth}@{}}
$\upsilon$ & \raggedright Heliozentr. Länge des Himmelskörpers in seiner Bahn.\tabularnewline
$r$ & \raggedright Wahrer Abstand von der Sonne.\tabularnewline
$\Delta$ & \raggedright Wahrer Abstand von der Erde.
\end{tabular}
\[
\left.\begin{array}{l}A\\B\\C\ \text{etc.}\end{array}\right\}\text{ Hilfswinkel.}
\]
\end{minipage}

\begin{center}I.\end{center}
\[
\begin{aligned}
1^{0}\quad\frac{\cos(\odot-\AscendingNode)\operatorname{tang}\ObservedBeta}{\sin(\odot-\alpha)}&=\operatorname{tang}A\\[5pt]
\frac{\sin A\operatorname{tang}(\odot-\AscendingNode)}{\sin(A+i)}&=\operatorname{tang}(\upsilon-\AscendingNode)
\end{aligned}
\]
\[
\begin{aligned}
2^{0}\quad\frac{\sin(\odot-\alpha)\operatorname{tang}i}{\cos(\odot-\AscendingNode)}&=\operatorname{tang}B\\[5pt]
\frac{\cos B\sin\ObservedBeta\operatorname{tang}(\odot-\AscendingNode)}{\sin(B+\ObservedBeta)\cos i}&=\operatorname{tang}(\upsilon-\AscendingNode)
\end{aligned}
\]
\[
\begin{aligned}
3^{0}\quad\frac{\sin(\odot-\AscendingNode)\operatorname{tang}\ObservedBeta}{\sin(\odot-\alpha)\operatorname{tang}i}&=\operatorname{tang}C\\[5pt]
\frac{\sin C\sin(\odot-\AscendingNode)}{\sin(C+\odot-\AscendingNode)\cos i}&=\operatorname{tang}(\upsilon-\AscendingNode)
\end{aligned}
\]
\[
\begin{aligned}
4^{0}\quad\frac{\cos(\odot-\AscendingNode)\operatorname{tang}\ObservedBeta}{\cos(\odot-\alpha)\operatorname{tang}i}&=\operatorname{tang}D\\[5pt]
\frac{\sin D\operatorname{tang}(\odot-\AscendingNode)\cos(\odot-\alpha)}{\sin(D+\odot-\alpha)\cos i}&=\operatorname{tang}(\upsilon-\AscendingNode)
\end{aligned}
\]
\textit{Anmerkung:} Da Winkel, die sich um $180^{0}$ unterscheiden, denselben Tangens
haben, ist noch eine Vorschrift dafür nötig, wie die Winkel $A$, $B$, $C$ usw.
und $\upsilon-\AscendingNode$ anhand ihrer Tangenswerte anzusetzen sind. Der Winkel
$\upsilon-\AscendingNode$ ist stets zwischen 0 und $180^{0}$ anzusetzen, wenn $\ObservedBeta$
\SourceLeaf{93}%
positiv (nördlich) ist. Ist die Breite dagegen südlich, so muss $\upsilon-\AscendingNode$
zwischen $180^{0}$ und $360^{0}$ oder, was dasselbe ist, zwischen $-180^{0}$ und 0 liegen.
Ist $\ObservedBeta=0$, so befindet sich der Himmelskörper an einem Knoten, und es
besteht nie Zweifel, ob es $\AscendingNode$ oder $\DescendingNode$ ist\HistoricalNote{3}.
Die Hilfswinkel $A$, $B$, $C$, $D$ und ebenso die folgenden $E$, $F$ usw. kann man
in dieser Hinsicht ganz nach Belieben ansetzen. Dabei sind die Vorzeichen $\pm$
selbstverständlich angemessen zu berücksichtigen. Im folgenden Beispiel habe ich
sie stets zwischen $-90^{0}$ und $+90^{0}$ gewählt. — Logarithmen, deren zugehörige
Größen negativ sind, habe ich durch ein dazugeschriebenes n gekennzeichnet.

\begin{center}II.\end{center}
\[
\begin{aligned}
5^{0}\quad\frac{\operatorname{tang}\ObservedBeta}{\sin(\alpha-\AscendingNode)}&=\operatorname{tang}E\\[5pt]
\frac{\sin E\sin(\odot-\AscendingNode)}{\sin(i-E)\sin(\upsilon-\AscendingNode)}&=\frac{r}{R}
\end{aligned}
\]
\[
\begin{aligned}
6^{0}\quad\operatorname{tang}i\sin(\alpha-\AscendingNode)&=\operatorname{tang}F\\[5pt]
\frac{\cos F\sin(\odot-\AscendingNode)\sin\ObservedBeta}{\sin(F-\ObservedBeta)\sin(\upsilon-\AscendingNode)\cos i}&=\frac{r}{R}
\end{aligned}
\]
\[
\begin{aligned}
7^{0}\quad\cos i\operatorname{tang}(\upsilon-\AscendingNode)&=\operatorname{tang}G\\[5pt]
\frac{\cos G\sin(\odot-\alpha)}{\sin(\alpha-\AscendingNode-G)\cos(\upsilon-\AscendingNode)}&=\frac{r}{R}
\end{aligned}
\]
\[
\begin{aligned}
8^{0}\quad\frac{\operatorname{tang}(\alpha-\AscendingNode)}{\cos i}&=\operatorname{tang}H\\[5pt]
\frac{\sin H\sin(\odot-\alpha)}{\sin(H-(\upsilon-\AscendingNode))\sin(\alpha-\AscendingNode)}&=\frac{r}{R}
\end{aligned}
\]
\[
\begin{aligned}
9^{0}\quad\frac{\operatorname{tang}\ObservedBeta}{\sin i\cos(\alpha-\AscendingNode)}&=\operatorname{tang}I\\[5pt]
\frac{\sin I\cos(\odot-\AscendingNode)}{\sin(\upsilon-\AscendingNode-I)}&=\frac{r}{R}
\end{aligned}
\]
\[
\begin{aligned}
10^{0}\quad\sin i\cos(\alpha-\AscendingNode)\operatorname{tang}(\upsilon-\AscendingNode)&=\operatorname{tang}K\\[5pt]
\frac{\cos K\sin\ObservedBeta\cos(\odot-\AscendingNode)}{\sin(K-\ObservedBeta)\cos(\upsilon-\AscendingNode)}&=\frac{r}{R}
\end{aligned}
\]
\[
\begin{aligned}
11^{0}\quad\frac{\sin C\sin(\odot-\alpha)}{\cos(C+\odot-\alpha)\operatorname{tang}(\odot-\AscendingNode)\cos i}&=\operatorname{tang}L\\[5pt]
\frac{\sin L}{\sin(\upsilon-\AscendingNode-L)\cos(\odot-\AscendingNode)}&=\frac{r}{R}
\end{aligned}
\]
\[
\begin{aligned}
12^{0}\quad\frac{\sin D\cos(\odot-\AscendingNode)}{\cos(D+\odot-\AscendingNode)\cos i}&=\operatorname{tang}M\\[5pt]
\frac{\sin M}{\sin(\upsilon-\AscendingNode-M)\cos(\odot-\AscendingNode)}&=\frac{r}{R}
\end{aligned}
\]
\begin{center}III.\end{center}
\[
13^{0}\quad\frac{r\sin(\upsilon-\AscendingNode)\sin i}{\sin\ObservedBeta}=\Delta
\]
\[
\begin{aligned}
14^{0}\quad\frac{R\sin E\sin(\odot-\AscendingNode)\sin i}{\sin(i-E)\sin\ObservedBeta}&\\[5pt]
&=\frac{R\cos E\sin(\odot-\AscendingNode)\sin i}{\sin(i-E)\sin(\alpha-\AscendingNode)\cos\ObservedBeta}\\[5pt]
&=\Delta
\end{aligned}
\]
\[
\begin{aligned}
15^{0}\quad\frac{R\cos F\sin(\odot-\AscendingNode)\operatorname{tang}i}{\sin(F-\ObservedBeta)}&\\[5pt]
&=\frac{R\sin F\sin(\odot-\AscendingNode)\sin(\alpha-\AscendingNode)}{\sin(F-\ObservedBeta)}\\[5pt]
&=\Delta
\end{aligned}
\]
Und so lassen sich noch weitere Ausdrücke für $\Delta$ ableiten, indem man
$13^{0}$ mit allen Formeln in II verbindet.

\par\Needspace{14\baselineskip}
\SourceLeaf{94}%
\begin{center}\textit{Beispiel.}\end{center}
\[
\begin{array}{rlrrr}
\AscendingNode&=&80^{0}&59'&12''07\\
\odot&=&281&1&34.99\\
\alpha&=&53&23&2.46\\
i&=&10&37&9.55
\end{array}
\]
\[
\begin{array}{@{}r@{\;}l@{\quad}l@{}}
\log\operatorname{tang}\ObservedBeta & {}=8.7349698\,\mathrm n & \ObservedBeta=-3^{0}\ 6'\ 33''561\quad\text{negativ oder südlich}\\
\log R & {}=9.9926158 &
\end{array}
\]
\[
\begin{array}{rlrrr}
\text{Folglich }\odot-\AscendingNode&=&200^{0}&2'&22''92\\
\odot-\alpha&=&227&38&32.53\\
\alpha-\AscendingNode&=&-27&36&9.61
\end{array}
\]
\par\medskip\noindent\begin{minipage}[t]{\textwidth}
\begin{center}$1^{0}$.\end{center}
{\smallskip\fontsize{8.7}{10.5}\selectfont
\noindent\begin{tabular*}{\linewidth}{@{}l@{\enspace}l@{\extracolsep{\fill}}p{1em}@{\extracolsep{0pt}}l@{\enspace}r@{\enspace}r@{\enspace}r@{\enspace}r@{}}
$\log\operatorname{tang}\ObservedBeta\enspace .\ .\ .\ .\ .$ & $8.7349698\,\mathrm n$ &  &  &  &  &  &\\
$\log\cos(\odot-\AscendingNode)\enspace .\ .\ .$ & $9.9728762\,\mathrm n$ &  &  &  &  &  &\\
$\text{Compl.}\ \log\sin(\odot-\alpha)\enspace .$ & $0.1313827\,\mathrm n$ &  &  &  &  &  &\\
\cline{1-2}
$\log\operatorname{tang}A\enspace .\ .\ .\ .\ .$ & $8.8392287\,\mathrm n$ &  & folglich & $A=-$ & $3^{0}$ & $57'$ & $2''136$\\
\cline{1-2}
$\log\sin A\enspace .\ .\ .\ .\ .$ & $8.8381955\,\mathrm n$ &  &  & $A+i=$ & $6$ & $40$ & $7.414$\\
$\log\operatorname{tang}(\odot-\AscendingNode)\enspace .\ .\ .$ & $9.5620014$ &  &  &  &  &  &\\
$\text{Compl.}\ \log\sin(A+i)\enspace .$ & $0.9350608$ &  &  &  &  &  &\\
\cline{1-2}
$\log\operatorname{tang}(\upsilon-\AscendingNode)\enspace .\ .\ .$ & $9.3352577\,\mathrm n$ &  & ferner & $\upsilon-\AscendingNode=-$ & $12^{0}$ & $12'$ & $37''942$\\
 &  &  & also & $\upsilon=.$ & $68$ & $46$ & $34.128$\\
\cline{4-8}
\end{tabular*}\par}
\end{minipage}\par
\par\medskip\noindent\begin{minipage}[t]{\textwidth}
\begin{center}$2^{0}$.\end{center}
{\smallskip\fontsize{8.7}{10.5}\selectfont
\noindent\begin{tabular*}{\linewidth}{@{}l@{\enspace}l@{\extracolsep{\fill}}p{1em}@{\extracolsep{0pt}}l@{\enspace}r@{\enspace}r@{\enspace}r@{\enspace}r@{}}
$\log\sin(\odot-\alpha)\enspace .\ .\ .\ .$ & $9.8686173\,\mathrm n$ &  &  &  &  &  &\\
$\log\operatorname{tang}i\enspace .\ .\ .\ .\ .$ & $9.2729872$ &  &  &  &  &  &\\
$\text{Compl.}\ \log\cos(\odot-\AscendingNode)\enspace .$ & $0.0271238\,\mathrm n$ &  &  &  &  &  &\\
\cline{1-2}
$\log\operatorname{tang}B\enspace .\ .\ .\ .\ .$ & $9.1687283$ &  & folglich & $B=$ & $8^{0}$ & $23'$ & $21''888$\\
\cline{1-2}
$\log\cos B\enspace .\ .\ .\ .\ .$ & $9.9953277$ &  &  & $B+\ObservedBeta=$ & $5$ & $16$ & $48.327$\\
$\log\sin\ObservedBeta\enspace .\ .\ .\ .\ .\ .$ & $8.7343300\,\mathrm n$ &  &  &  &  &  &\\
$\log\operatorname{tang}(\odot-\AscendingNode)\enspace .\ .\ .$ & $9.5620014$ &  &  &  &  &  &\\
$\text{Compl.}\ \log\sin(B+\ObservedBeta)\enspace .$ & $1.0360961$ &  &  &  &  &  &\\
$\text{Compl.}\ \log\cos i\enspace .\ .\ .$ & $0.0075025$ &  &  &  &  &  &\\
\cline{1-2}
$\log\operatorname{tang}(\upsilon-\AscendingNode)\enspace .\ .\ .$ & $9.3352577\,\mathrm n$ &  & \multicolumn{5}{@{}l@{}}{wie oben}\\
\end{tabular*}\par}
\end{minipage}\par

\SourceLeaf{95}%

\par\medskip\noindent\begin{minipage}[t]{\textwidth}
\begin{center}$3^{0}$.\end{center}
{\smallskip\fontsize{8.7}{10.5}\selectfont
\noindent\begin{tabular*}{\linewidth}{@{}l@{\enspace}l@{\extracolsep{\fill}}p{1em}@{\extracolsep{0pt}}l@{\enspace}r@{\enspace}r@{\enspace}r@{\enspace}r@{}}
$\log\sin(\odot-\AscendingNode)\enspace .\ .\ .$ & $9.5348776\,\mathrm n$ &  &  &  &  &  &\\
$\log\operatorname{tang}\ObservedBeta\enspace .\ .\ .\ .\ .$ & $8.7349698\,\mathrm n$ &  &  &  &  &  &\\
$\text{Compl.}\ \log\sin(\odot-\alpha)\enspace .$ & $0.1313827\,\mathrm n$ &  &  &  &  &  &\\
$\text{Compl.}\ \log\operatorname{tang}i\enspace .\ .\ .$ & $0.7270128$ &  &  &  &  &  &\\
\cline{1-2}
$\log\operatorname{tang}C\enspace .\ .\ .\ .\ .$ & $9.1282429\,\mathrm n$ &  & also & $C=-$ & $7^{0}$ & $39'$ & $7''056$\\
\cline{1-2}
$\log\sin C\enspace .\ .\ .\ .\ .$ & $9.1243583\,\mathrm n$ &  &  & $C+\odot-\AscendingNode=$ & $192$ & $23$ & $15.864$\\
$\log\sin(\odot-\AscendingNode)\enspace .\ .\ .$ & $9.5348776\,\mathrm n$ &  &  &  &  &  &\\
$\text{Compl.}\ \log\sin(C+\odot-\AscendingNode)$ & $0.6685194\,\mathrm n$ &  &  &  &  &  &\\
$\text{Compl.}\ \log\cos i\enspace .\ .\ .$ & $0.0075025$ &  &  &  &  &  &\\
\cline{1-2}
$\log\operatorname{tang}(\upsilon-\AscendingNode)\enspace .\ .\ .$ & $9.3352578\,\mathrm n$ &  & \multicolumn{5}{@{}l@{}}{wie vorhin.}\\
\end{tabular*}\par}
\end{minipage}\par
\par\medskip\noindent\begin{minipage}[t]{\textwidth}
\begin{center}$4^{0}$.\end{center}
{\smallskip\fontsize{8.7}{10.5}\selectfont
\noindent\begin{tabular*}{\linewidth}{@{}l@{\enspace}l@{\extracolsep{\fill}}p{1em}@{\extracolsep{0pt}}l@{\enspace}r@{\enspace}r@{\enspace}r@{\enspace}r@{}}
$\log\cos(\odot-\AscendingNode)\enspace .\ .\ .$ & $9.9728762\,\mathrm n$ &  &  &  &  &  &\\
$\log\operatorname{tang}\ObservedBeta\enspace .\ .\ .\ .\ .$ & $8.7349698\,\mathrm n$ &  &  &  &  &  &\\
$\text{Compl.}\ \log\cos(\odot-\alpha)\enspace .$ & $0.1714973\,\mathrm n$ &  &  &  &  &  &\\
$\text{Compl.}\ \log\operatorname{tang}i\enspace .\ .\ .$ & $0.7270128$ &  &  &  &  &  &\\
\cline{1-2}
$\log\operatorname{tang}D\enspace .\ .\ .\ .\ .$ & $9.6063561\,\mathrm n$ &  & also & $D=-$ & $21^{0}$ & $59'$ & $51''182$\\
\cline{1-2}
$\log\sin D\enspace .\ .\ .\ .\ .$ & $9.5735295\,\mathrm n$ &  &  & $D+\odot-\alpha=$ & $205$ & $38$ & $41.348$\\
$\log\operatorname{tang}(\odot-\AscendingNode)\enspace .\ .\ .$ & $9.5620014$ &  &  &  &  &  &\\
$\log\cos(\odot-\alpha)\enspace .\ .\ .\ .$ & $9.8285027\,\mathrm n$ &  &  &  &  &  &\\
$\text{Compl.}\ \log\sin(D+\odot-\alpha)$ & $0.3637217\,\mathrm n$ &  &  &  &  &  &\\
$\text{Compl.}\ \log\cos i\enspace .\ .\ .$ & $0.0075025$ &  &  &  &  &  &\\
\cline{1-2}
$\log\operatorname{tang}(\upsilon-\AscendingNode)\enspace .\ .\ .\ .$ & $9.3352578\,\mathrm n$ &  & \multicolumn{5}{@{}l@{}}{wie oben.}\\
\end{tabular*}\par}
\end{minipage}\par
\par\medskip\noindent\begin{minipage}[t]{\textwidth}
\begin{center}$5^{0}$.\end{center}
{\smallskip\fontsize{8.7}{10.5}\selectfont
\noindent\begin{tabular*}{\linewidth}{@{}l@{\enspace}l@{\extracolsep{\fill}}p{1em}@{\extracolsep{0pt}}l@{\enspace}r@{\enspace}r@{\enspace}r@{\enspace}r@{}}
$\log\operatorname{tang}\ObservedBeta\enspace .\ .\ .\ .\ .$ & $8.7349698\,\mathrm n$ &  &  &  &  &  &\\
$\log\sin(\alpha-\AscendingNode)\enspace .\ .\ .\ .$ & $9.6658973\,\mathrm n$ &  &  &  &  &  &\\
\cline{1-2}
$\log\operatorname{tang}E\enspace .\ .\ .\ .\ .$ & $9.0690725$ &  & also & $E=$ & $6^{0}$ & $41'$ & $12''412$\\
\cline{1-2}
$\log\sin E\enspace .\ .\ .\ .\ .$ & $9.0661081$ &  &  & $i-E=$ & $3$ & $55$ & $57.138$\\
$\log\sin(\odot-\AscendingNode)\enspace .\ .\ .\ .$ & $9.5348776\,\mathrm n$ &  &  &  &  &  &\\
$\text{Compl.}\ \log\sin(i-E)\enspace .$ & $1.1637907$ &  &  &  &  &  &\\
$\text{Compl.}\ \log\sin(\upsilon-\AscendingNode)\enspace .$ & $0.6746802\,\mathrm n$ &  & \multicolumn{5}{@{}c@{}}{ferner}\\
\cline{1-2}
$\log\frac{r}{R}\enspace .\ .\ .\ .\ .\ .\ .$ & $0.4394566$ &  & \multicolumn{5}{@{}l@{}}{$\log r=\log R+\log\dfrac{r}{R}=0.4320724$}\\
\end{tabular*}\par}
\end{minipage}\par

\SourceLeaf{96}%

\par\medskip\noindent\begin{minipage}[t]{\textwidth}
\begin{center}$6^{0}$.\end{center}
{\smallskip\fontsize{8.7}{10.5}\selectfont
\noindent\begin{tabular*}{\linewidth}{@{}l@{\enspace}l@{\extracolsep{\fill}}p{1em}@{\extracolsep{0pt}}l@{\enspace}r@{\enspace}r@{\enspace}r@{\enspace}r@{}}
$\log\operatorname{tang}i\enspace .\ .\ .\ .\ .$ & $9.2729872$ &  &  &  &  &  &\\
$\log\sin(\alpha-\AscendingNode)\enspace .\ .\ .\ .$ & $9.6658973\,\mathrm n$ &  &  &  &  &  &\\
\cline{1-2}
$\log\operatorname{tang}F\enspace .\ .\ .\ .\ .$ & $8.9388845\,\mathrm n$ &  & daher & $F=-$ & $4^{0}$ & $57'$ & $53''955$\\
\cline{1-2}
$\log\cos F\enspace .\ .\ .\ .\ .$ & $9.9983674$ &  &  & $F-\ObservedBeta=-$ & $1$ & $51$ & $20.394$\\
$\log\sin\ObservedBeta\enspace .\ .\ .\ .\ .\ .$ & $8.7343300\,\mathrm n$ &  &  &  &  &  &\\
$\log\sin(\odot-\AscendingNode)\enspace .\ .\ .\ .$ & $9.5348776\,\mathrm n$ &  &  &  &  &  &\\
$\text{Compl.}\ \log\sin(F-\ObservedBeta)\enspace .$ & $1.4896990\,\mathrm n$ &  &  &  &  &  &\\
$\text{Compl.}\ \log\sin(\upsilon-\AscendingNode)\enspace .$ & $0.6746802\,\mathrm n$ &  &  &  &  &  &\\
$\text{Compl.}\ \log\cos i\enspace .\ .\ .$ & $0.0075025$ &  &  &  &  &  &\\
\cline{1-2}
$\log\frac{r}{R}\enspace .\ .\ .\ .\ .\ .\ .$ & $0.4394567$ &  & \multicolumn{5}{@{}l@{}}{annähernd wie zuvor.}\\
\end{tabular*}\par}
\end{minipage}\par
\par\medskip\noindent\begin{minipage}[t]{\textwidth}
\begin{center}$7^{0}$.\end{center}
{\smallskip\fontsize{8.7}{10.5}\selectfont
\noindent\begin{tabular*}{\linewidth}{@{}l@{\enspace}l@{\extracolsep{\fill}}p{1em}@{\extracolsep{0pt}}l@{\enspace}r@{\enspace}r@{\enspace}r@{\enspace}r@{}}
$\log\cos i\enspace .\ .\ .\ .\ .\ .$ & $9.9924975$ &  &  &  &  &  &\\
$\log\operatorname{tang}(\upsilon-\AscendingNode)\enspace .\ .\ .$ & $9.3352577\,\mathrm n$ &  &  &  &  &  &\\
\cline{1-2}
$\log\operatorname{tang}G\enspace .\ .\ .\ .\ .$ & $9.3277552\,\mathrm n$ &  & also & $G=-$ & $12^{0}$ & $0'$ & $27''118$\\
\cline{1-2}
$\log\cos G\enspace .\ .\ ,\ .\ .$ & $9.9903922$ &  &  & $\alpha-\AscendingNode-G=-$ & $15$ & $35$ & $42.492$\\
$\log\sin(\odot-\alpha)\enspace .\ .\ .\ .$ & $9.8686173\,\mathrm n$ &  &  &  &  &  &\\
$\text{Compl.}\ \log\sin(\alpha-\AscendingNode-G)$ & $0.5705092\,\mathrm n$ &  &  &  &  &  &\\
$\text{Compl.}\ \log\cos(\upsilon-\AscendingNode)\enspace .$ & $0.0099379$ &  &  &  &  &  &\\
\cline{1-2}
$\log\frac{r}{R}\enspace .\ .\ .\ .\ .\ .\ .$ & $0.4394566$ &  & \multicolumn{5}{@{}l@{}}{wie oben.}\\
\end{tabular*}\par}
\end{minipage}\par
\par\medskip\noindent\begin{minipage}[t]{\textwidth}
\begin{center}$8^{0}$.\end{center}
{\smallskip\fontsize{8.7}{10.5}\selectfont
\noindent\begin{tabular*}{\linewidth}{@{}l@{\enspace}l@{\extracolsep{\fill}}p{1em}@{\extracolsep{0pt}}l@{\enspace}r@{\enspace}r@{\enspace}r@{\enspace}r@{}}
$\log\operatorname{tang}(\alpha-\AscendingNode)\enspace .\ .\ .$ & $9.7183744\,\mathrm n$ &  &  &  &  &  &\\
$\log\cos i\enspace .\ .\ .\ .\ .\ .$ & $9.9924975$ &  &  &  &  &  &\\
\cline{1-2}
$\log\operatorname{tang}H\enspace .\ .\ .\ .\ .$ & $9.7258769\,\mathrm n$ &  & folglich & $H=-$ & $28^{0}$ & $0'$ & $39''879$\\
\cline{1-2}
$\log\sin H\enspace .\ .\ .\ .\ .$ & $9.6717672\,\mathrm n$ &  &  & $H-(\upsilon-\AscendingNode)=-$ & $15$ & $48$ & $1.937$\\
$\log\sin(\odot-\alpha)\enspace .\ .\ .\ .$ & $9.8686173\,\mathrm n$ &  &  &  &  &  &\\
$\text{Compl.}\ \log\sin(H-(\upsilon-\AscendingNode))$ & $0.5649695\,\mathrm n$ &  &  &  &  &  &\\
$\text{Compl.}\ \log\sin(\alpha-\AscendingNode)\enspace .$ & $0.3341027\,\mathrm n$ &  &  &  &  &  &\\
\cline{1-2}
$\log\frac{r}{R}\enspace .\ .\ .\ .\ .\ .\ .$ & $0.4394567$ &  & \multicolumn{5}{@{}l@{}}{wie vorher.}\\
\end{tabular*}\par}
\end{minipage}\par

\SourceLeaf{97}%

\par\medskip\noindent\begin{minipage}[t]{\textwidth}
\begin{center}$9^{0}$.\end{center}
{\smallskip\fontsize{8.7}{10.5}\selectfont
\noindent\begin{tabular*}{\linewidth}{@{}l@{\enspace}l@{\extracolsep{\fill}}p{1em}@{\extracolsep{0pt}}l@{\enspace}r@{\enspace}r@{\enspace}r@{\enspace}r@{}}
$\log\operatorname{tang}\ObservedBeta\enspace .\ .\ .\ .\ .$ & $8.7349698\,\mathrm n$ &  &  &  &  &  &\\
$\text{Compl.}\ \log\sin i\enspace .\ .\ .$ & $0.7345153$ &  &  &  &  &  &\\
$\text{Compl.}\ \log\cos(\alpha-\AscendingNode)\enspace .$ & $0.0524771$ &  &  &  &  &  &\\
\cline{1-2}
$\log\operatorname{tang}I\enspace .\ .\ .\ .\ .$ & $9.5219622\,\mathrm n$ &  & hieraus & $I=-$ & $18^{0}$ & $23'$ & $55''334$\\
\cline{1-2}
$\log\sin I\enspace .\ .\ .\ .\ .\ .$ & $9.4991749\,\mathrm n$ &  &  & $\upsilon-\AscendingNode-I=$ & $6$ & $11$ & $17.392$\\
$\log\cos(\odot-\AscendingNode)\enspace .\ .\ .$ & $9.9728762\,\mathrm n$ &  &  &  &  &  &\\
$\text{Compl.}\ \log\sin(\upsilon-\AscendingNode-I)$ & $0.9674054$ &  &  &  &  &  &\\
\cline{1-2}
$\log\frac{r}{R}\enspace .\ .\ .\ .\ .\ .\ .$ & $0.4394565$ &  & \multicolumn{5}{@{}l@{}}{wie vorhin.}\\
\end{tabular*}\par}
\end{minipage}\par
\par\medskip\noindent\begin{minipage}[t]{\textwidth}
\begin{center}$10^{0}$.\end{center}
\begin{center}In der Nähe des Knotens weniger genau.\end{center}
{\smallskip\fontsize{8.7}{10.5}\selectfont
\noindent\begin{tabular*}{\linewidth}{@{}l@{\enspace}l@{\extracolsep{\fill}}p{1em}@{\extracolsep{0pt}}l@{\enspace}r@{\enspace}r@{\enspace}r@{\enspace}r@{}}
$\log\sin i\enspace .\ .\ .\ .\ .\ .$ & $9.2654847$ &  &  &  &  &  &\\
$\log\cos(\alpha-\AscendingNode)\enspace .\ .\ .\ .$ & $9.9475229$ &  &  &  &  &  &\\
$\log\operatorname{tang}(\upsilon-\AscendingNode)\enspace .\ .\ .\ .$ & $9.3352577\,\mathrm n$ &  &  &  &  &  &\\
\cline{1-2}
$\log\operatorname{tang}K\enspace .\ .\ .\ .\ .$ & $8.5482653\,\mathrm n$ &  & also & $K=-$ & $2^{0}$ & $1'$ & $26''344$\\
\cline{1-2}
$\log\cos K\enspace .\ .\ .\ .\ .$ & $9.9997290$ &  &  & $K-\ObservedBeta=$ & $1$ & $5$ & $7.217$\\
$\log\sin\ObservedBeta\enspace .\ .\ .\ .\ .\ .$ & $8.7343300\,\mathrm n$ &  &  &  &  &  &\\
$\log\cos(\odot-\AscendingNode)\enspace .\ .\ .$ & $9.9728762\,\mathrm n$ &  &  &  &  &  &\\
$\text{Compl.}\ \log\sin(K-\ObservedBeta)\enspace .$ & $1.7225836$ &  &  &  &  &  &\\
$\text{Compl.}\ \log\cos(\upsilon-\AscendingNode)\enspace .$ & $0.0099379$ &  &  &  &  &  &\\
\cline{1-2}
$\log\frac{r}{R}\enspace .\ .\ .\ .\ .\ .\ .$ & $0.4394567$ &  & \multicolumn{5}{@{}l@{}}{wie vorhin.}\\
\end{tabular*}\par}
\end{minipage}\par
\par\medskip\noindent\begin{minipage}[t]{\textwidth}
\begin{center}$11^{0}$.\end{center}
\[C+\odot-\alpha=219^{0}\ 59'\ 25''474\]
{\smallskip\fontsize{8.7}{10.5}\selectfont
\noindent\begin{tabular*}{\linewidth}{@{}l@{\enspace}l@{\extracolsep{\fill}}p{1em}@{\extracolsep{0pt}}l@{\enspace}r@{\enspace}r@{\enspace}r@{\enspace}r@{}}
$\log\sin C\enspace .\ .\ .\ .\ .$ & $9.1243583\,\mathrm n$ &  &  &  &  &  &\\
$\log\sin(\odot-\alpha)\enspace .\ .\ .\ .$ & $9.8686173\,\mathrm n$ &  &  &  &  &  &\\
$\text{Compl.}\ \log\cos(C+\odot-\alpha)$ & $0.1156850\,\mathrm n$ &  &  &  &  &  &\\
$\text{Compl.}\ \log\operatorname{tang}(\odot-\AscendingNode)$ & $0.4379986$ &  &  &  &  &  &\\
$\text{Compl.}\ \log\cos i\enspace .\ .\ .$ & $0.0075025$ &  &  &  &  &  &\\
\cline{1-2}
$\log\operatorname{tang}L\enspace .\ .\ .\ .\ .$ & $9.5541617\,\mathrm n$ &  & also & $L=-$ & $19^{0}$ & $42'$ & $32''533$\\
\cline{1-2}
$\log\sin L\enspace .\ .\ .\ .\ .\ .$ & $9.5279439\,\mathrm n$ &  &  & $\upsilon-\AscendingNode-L=$ & $7$ & $29$ & $54.591$\\
$\text{Compl.}\ \log\sin(\upsilon-\AscendingNode-L)$ & $0.8843888$ &  &  &  &  &  &\\
$\text{Compl.}\ \log\cos(\odot-\AscendingNode)\enspace .$ & $0.0271238\,\mathrm n$ &  &  &  &  &  &\\
\cline{1-2}
$\log\frac{r}{R}\enspace .\ .\ .\ .\ .\ .\ .$ & $0.4394565$ &  & \multicolumn{5}{@{}l@{}}{wie zuvor.}\\
\end{tabular*}\par}
\end{minipage}\par

\SourceLeaf{98}%

\par\medskip\noindent\begin{minipage}[t]{\textwidth}
\begin{center}$12^{0}$.\end{center}
\[D+\odot-\AscendingNode=178^{0}\ 2'\ 31''738\]
{\smallskip\fontsize{8.7}{10.5}\selectfont
\noindent\begin{tabular*}{\linewidth}{@{}l@{\enspace}l@{\extracolsep{\fill}}p{1em}@{\extracolsep{0pt}}l@{\enspace}r@{\enspace}r@{\enspace}r@{\enspace}r@{}}
$\log\sin D\enspace .\ .\ .\ .\ .$ & $9.5735295\,\mathrm n$ &  &  &  &  &  &\\
$\log\cos(\odot-\AscendingNode)\enspace .\ .\ .$ & $9.9728762\,\mathrm n$ &  &  &  &  &  &\\
$\text{Compl.}\ \log\cos(D+\odot-\AscendingNode)$ & $0.0002536\,\mathrm n$ &  &  &  &  &  &\\
$\text{Compl.}\ \log\cos i\enspace .\ .\ .$ & $0.0075025$ &  &  &  &  &  &\\
\cline{1-2}
$\log\operatorname{tang}(M=L)\enspace .\ .\ .$ & $9.5541618\,\mathrm n$ &  & \multicolumn{5}{@{}l@{}}{wie oben in $11^{0}$.}\\
\end{tabular*}\par}

Der übrige Teil der Rechnung verläuft genauso wie dort.


\end{minipage}\par
\par\medskip\noindent\begin{minipage}[t]{\textwidth}
\begin{center}$13^{0}$.\end{center}
\begin{center}\small
\begin{tabular}{@{}lcl@{}}
$\log r$ & $.\ .\ .\ .\ .\ .\ .$ & $0.4320724$\\
$\log\sin(\upsilon-\AscendingNode)$ & $.\ .\ .\ .$ & $9.3253198\,\mathrm n$\\
$\log\sin i$ & $.\ .\ .\ .\ .\ .$ & $9.2654847$\\
$\text{Compl.}\ \log\sin\ObservedBeta$ & $.\ .\ .$ & $1.2656700\,\mathrm n$\\
\hline
$\log\Delta$ & $.\ .\ .\ .\ .\ .\ .$ & $0.2885469$\\
\end{tabular}
\end{center}
\par\medskip\noindent\rule{.3\textwidth}{.3pt}

\end{minipage}\par
